\casechapter{Completing the Integral}{Lebesgue and the escaping limits}{Completion}{ch:lebesgue}
\begin{refsection}[chapters/01-lebesgue.bib]
\newcommand{\lebRiem}{\mathcal{R}}
\newcommand{\lebsymdiff}{\mathbin{\triangle}}

\begin{chaptersummary}
Riemann's integral is defined by a limit, yet it is not closed under limits: pointwise limits escape it, sequences that converge in the mean have nothing to converge to, and a bounded derivative need not be integrable. Lebesgue measured the sets on which a function takes its values and summed over the range instead of the domain. From a distance the move is a completion: $L^1$ is the completion of the step functions, and Lebesgue's integral is the only continuous extension of the elementary one. In the completed space the limit theorems become statements about continuity: Arzelà's bounded convergence theorem shrinks from a page of compactness arguments to two lines, Fourier series become coordinates, and every bounded derivative can be integrated.
\end{chaptersummary}

\section{Before: three leaks in Riemann's integral}\label{sec:lebesgue:before}

Riemann defined his integral in 1854, in the Habilitationsschrift on trigonometric series that was published only after his death \parencite{Riemann1867}.\sidenote{The integral appears in a paper about trigonometric series because Fourier coefficients are integrals: before asking which functions have Fourier series, one must ask which functions can be integrated.} Darboux later recast it with upper and lower sums, and in that form it reached every calculus course. We use an equivalent formulation with step functions, because step functions will also be the raw material of the move.

\begin{definition}[Step functions and the Riemann integral]\label[definition]{def:lebesgue:riemann}
A \term{step function} on $[a,b]$ is a finite linear combination $s=\sum_i c_i\ind_{I_i}$, with real coefficients, of indicator functions of subintervals $I_i\subseteq[a,b]$; its \term{elementary integral} is $\int s=\sum_i c_i\abs{I_i}$, where $\abs{I}$ is the length of $I$. If $a=x_0<x_1<\dots<x_m=b$ are the points $a$, $b$, and the endpoints of the $I_i$, then $s$ is constant on each $(x_{j-1},x_j)$; these are its \term{open intervals of constancy}. A function $f\colon[a,b]\to\R$ is \term{Riemann integrable} if for every $\eps>0$ there are step functions $\varphi\le f\le\psi$ with $\int(\psi-\varphi)<\eps$; its integral $\int_a^b f$ is then the common value of $\sup_\varphi\int\varphi$ and $\inf_\psi\int\psi$. We write $\lebRiem[a,b]$ for the class of Riemann integrable functions.
\end{definition}

The elementary integral does not depend on how $s$ is written, and changing $s$ at finitely many points does not change it. Riemann integrable functions are bounded, the class $\lebRiem[a,b]$ is closed under sums, products, and absolute values, and the integral is linear and monotone on it. Darboux's upper and lower sums are the integrals of the step functions $\psi$ and $\varphi$ built from a partition, so the definition agrees with the familiar one.\index{Riemann integral}\index{step function}

\subsection{Limits escape}\label{sec:lebesgue:escape}

Enumerate the rationals of $[0,1]$ as $q_1,q_2,\dots$, and let $f_n$ be the indicator function of the finite set $\{q_1,\dots,q_n\}$. Each $f_n$ is a step function with integral $0$ (a point is a degenerate interval), and the sequence increases at every point to \term{Dirichlet's function} $D=\ind_{\Q\cap[0,1]}$.\index{Dirichlet's function} The limit is not Riemann integrable. A step function $\psi\ge D$ is at least $1$ on each of its open intervals of constancy, since such an interval contains a rational, and a step function $\varphi\le D$ is at most $0$ on each of them, since such an interval contains an irrational; hence $\int(\psi-\varphi)\ge1$ for every such pair. An increasing sequence of integrable functions, bounded by $1$, with integrals all equal to $0$, has escaped the class.

The classical cure is uniform convergence: a uniform limit of Riemann integrable functions is Riemann integrable, with the limit of the integrals as its integral. Most situations provide less, and the textbooks supplemented uniform convergence with arguments tailored to each case, such as this one for series.

\begin{proposition}[Term-by-term integration, Riemann's version]\label[proposition]{prop:lebesgue:dini}
Let $f_1,f_2,\dots$ be nonnegative continuous functions on $[a,b]$ whose series converges at every point to a continuous function $f$. Then $\int_a^b f=\sum_n\int_a^b f_n$.
\end{proposition}

\begin{proof}
The plan is to upgrade pointwise convergence to uniform convergence by compactness, an argument known as Dini's theorem.\index{Dini's theorem} The remainders $g_N=f-\sum_{n\le N}f_n$ are continuous, and they decrease to $0$ at every point because the $f_n$ are nonnegative. Fix $\eps>0$. The sets $K_N=\{x\in[a,b]:g_N(x)\ge\eps\}$ are closed because $g_N$ is continuous, they decrease with $N$ because $g_N$ does, and their intersection is empty because $g_N(x)\to0$ for each $x$. A decreasing sequence of nonempty closed subsets of the compact interval $[a,b]$ has nonempty intersection, so some $K_N$ is empty. Then $0\le g_M<\eps$ on $[a,b]$ for every $M\ge N$, and integrating gives $0\le\int_a^b f-\sum_{n\le M}\int_a^b f_n\le\eps(b-a)$.
\end{proof}

Continuity of the $f_n$ and of $f$ makes the sets $K_N$ closed, and compactness does the rest; without continuity the argument collapses. If the $f_n$ are merely Riemann integrable and so is their sum $f$, then the conclusion survives, because the partial sums lie between $0$ and $\sup f$; but the proof now needs Arzelà's theorem.

\begin{theorem}[\textcite{Arzela1885}]\label{thm:lebesgue:arzela}
Let $f_n\in\lebRiem[a,b]$ satisfy $\abs{f_n}\le M$ for all $n$, where $M$ is a constant, and suppose that $f_n(x)\to f(x)$ for every $x\in[a,b]$, where $f\in\lebRiem[a,b]$. Then $\int_a^b f_n\to\int_a^b f$.
\end{theorem}
\index{Arzelà's bounded convergence theorem}

The hypothesis $f\in\lebRiem[a,b]$ cannot be dropped, by Dirichlet's example, and in applications it can rarely be checked without knowing the answer. The theorem is true, and proofs that avoid measure theory exist \parencite{Luxemburg1971,Lewin1986}, but they are delicate.\sidenote{Osgood had treated the special case in which the $f_n$ and $f$ are continuous, in a memoir of 1897.} The aside gives Lewin's short proof in full. A reader who skims it should notice where the difficulty lies.

\begin{aside}[An elementary proof of Arzelà's theorem]
Call a finite, possibly empty, union of closed intervals an \emph{elementary set}, and let $\abs{E}$ be its length, the elementary integral of $\ind_E$. Unions and intersections of elementary sets are elementary, and $\abs{E\cup F}+\abs{E\cap F}=\abs{E}+\abs{F}$ because $\ind_{E\cup F}+\ind_{E\cap F}=\ind_E+\ind_F$. For an arbitrary set $A\subseteq[a,b]$, let $\alpha(A)$ be the supremum of $\abs{E}$ over elementary sets $E\subseteq A$.

\emph{Lemma.} If $A_1\supseteq A_2\supseteq\cdots$ are subsets of $[a,b]$ with empty intersection, then $\alpha(A_n)\to0$.

\emph{Proof of the lemma.} The numbers $\alpha(A_n)$ decrease; suppose that $\alpha(A_n)\ge\delta>0$ for all $n$. Choose elementary sets $E_n\subseteq A_n$ with $\abs{E_n}>\alpha(A_n)-\delta2^{-n}$, and put $H_n=E_1\cap\dots\cap E_n$. The set $H_n\cup E_{n+1}$ is an elementary subset of $A_n$, so its length is at most $\alpha(A_n)$, and inclusion--exclusion gives $\abs{H_{n+1}}\ge\abs{H_n}+\abs{E_{n+1}}-\alpha(A_n)$. By induction, $\abs{H_n}>\alpha(A_n)-\delta(1-2^{-n})\ge\delta2^{-n}$. So the $H_n$ form a decreasing sequence of nonempty compact sets; their intersection is nonempty, and it lies in every $A_n$, a contradiction.

\emph{Proof of the theorem.} The functions $g_n=\abs{f_n-f}$ are Riemann integrable, bounded by $2M$, and tend to $0$ pointwise, and $\abs{\int f_n-\int f}\le\int g_n$; so it suffices to show that $\int g_n\to0$. Fix $\eps>0$ and let $A_n$ be the set of $x$ with $g_k(x)\ge\eps$ for some $k\ge n$; these sets decrease and have empty intersection. Let $\varphi\le g_n$ be a step function. The open intervals of constancy on which $\varphi\ge\eps$ lie in $A_n$, so their total length is at most $\alpha(A_n)$ (shrink them slightly to closed intervals); on the other intervals, $\varphi<\eps$. Since $\varphi\le2M$, we get $\int\varphi\le2M\alpha(A_n)+\eps(b-a)$, and the supremum over $\varphi$ gives $\int g_n\le2M\alpha(A_n)+\eps(b-a)$. By the lemma, $\limsup_n\int g_n\le\eps(b-a)$, and $\eps$ was arbitrary.

\emph{Where the difficulty lies.} The sets $A_n$ are not finite unions of intervals, and Riemann's theory gives them no length. The lemma manufactures just enough of a length, an inner content, to prove one continuity property, and it pays with a compactness argument and an induction. The next theorem of this kind must pay again.
\end{aside}

\subsection{Spaces are incomplete}\label{sec:lebesgue:incomplete}

Measure the distance between two integrable functions by their mean deviation $\norm{f-g}_1=\int_a^b\abs{f-g}$, a seminorm, since functions that differ at finitely many points are at distance zero. A sequence $(f_n)$ is \term{Cauchy} for this distance if $\norm{f_n-f_m}_1\to0$ as $n,m\to\infty$. Such sequences arise constantly, as partial sums of orthogonal expansions and as approximations of a shape by unions of boxes, and one expects them to converge. In Riemann's class they need not.

Start from $C_0=[0,1]$. At stage $n\ge1$ the set $C_{n-1}$ is a union of $2^{n-1}$ disjoint closed intervals; remove from the middle of each an open interval of length $4^{-n}$.\marginaside{\begin{tikzpicture}[x=3.9cm,y=1cm]
  \foreach \row/\list in {0/{0/1},1/{0/0.375,0.625/1},2/{0/0.15625,0.21875/0.375,0.625/0.78125,0.84375/1},3/{0/0.0703,0.0859/0.15625,0.21875/0.2891,0.3047/0.375,0.625/0.6953,0.7109/0.78125,0.84375/0.9141,0.9297/1}}{
    \foreach \s/\e in \list {\fill[sea] (\s,-0.32*\row) rectangle (\e,-0.32*\row+0.11);}
  }
\end{tikzpicture}\par\smallskip The fat Cantor set: stages $C_0,\dots,C_3$, of lengths $1,\frac34,\frac58,\frac9{16}$.} The length removed at stage $n$ is $2^{n-1}4^{-n}=2^{-n-1}$, so $C_n$ has total length $\frac12+2^{-n-1}$, and each of its $2^n$ intervals has length $2^{-n-1}+2^{-2n-1}$. The \term{fat Cantor set} $C=\bigcap_n C_n$ is closed and contains no nondegenerate interval, yet, in a language we do not have yet, it has length $\frac12$.\index{fat Cantor set}\sidenote{Nowhere dense sets of positive length appear in work of H.~J.~S. Smith (1875) and of Volterra (1881); the name Smith--Volterra--Cantor set is modern.}

\begin{proposition}[Riemann's class is incomplete]\label[proposition]{prop:lebesgue:incomplete}
The step functions $s_n=\ind_{C_n}$ form a Cauchy sequence for $\norm{\cdot}_1$, but no $g\in\lebRiem[0,1]$ satisfies $\norm{s_n-g}_1\to0$.
\end{proposition}

\begin{proof}
For $m>n$ we have $C_m\subseteq C_n$, so $\norm{s_n-s_m}_1=\abs{C_n}-\abs{C_m}<2^{-n-1}$, and the sequence is Cauchy. Suppose, for a contradiction, that $g\in\lebRiem[0,1]$ and $\norm{s_n-g}_1\to0$. Two facts follow. First, $\int g=\lim_n\int s_n=\frac12$, because $\abs{\int s_n-\int g}\le\norm{s_n-g}_1$. Second, $\int_I\abs{g}=0$ for every open interval $I$ removed in the construction: if $I$ was removed at stage $k$, then $s_n=0$ on $I$ for $n\ge k$, so $\int_I\abs{g}=\int_I\abs{s_n-g}\le\norm{s_n-g}_1\to0$.

Now let $\varphi\le g$ be a step function, and let $J$ be one of its open intervals of constancy, with value $c$. Since $C$ contains no nondegenerate interval, $J$ meets some removed interval $I$, and $J\cap I$ is a nondegenerate interval on which $c\le g\le\abs{g}$. If $c$ were positive, then $\int_{J\cap I}\abs{g}\ge c\abs{J\cap I}>0$, against the second fact. So every value of $\varphi$ on its intervals of constancy is at most $0$, and $\int\varphi\le0$. The supremum over $\varphi$ gives $\int g\le0$, against the first fact.
\end{proof}

The escaping limit is easy to name, since $s_n$ decreases to $\ind_C$, but by the proposition no Riemann integrable function is the limit of the $s_n$; \cref{ex:lebesgue:cantorlimit} will show that none even agrees with $\ind_C$ almost everywhere. The hole matters for Fourier series. By Bessel's inequality the Fourier coefficients of a Riemann integrable function are square summable, and the natural converse, that every square-summable sequence is the coefficient sequence of some function, fails in Riemann's class: the limits of the Fourier coefficients of $s_n$ form a square-summable sequence that belongs to no Riemann integrable function (Exercise~\ref{ex:lebesgue:fourier}). The theorem that ought to hold has no space in which to hold.

\subsection{Derivatives escape}\label{sec:lebesgue:derivatives}

The fundamental theorem of calculus, $\int_a^bF'=F(b)-F(a)$, carries a hypothesis in Riemann's theory: the derivative must be Riemann integrable. Derivatives have the intermediate value property, as Darboux showed, so they cannot jump, and one might hope that bounded derivatives are always integrable. \textcite{Volterra1881} showed otherwise.\index{Volterra's function} On each open interval removed in the construction of a fat Cantor set $C$, place a copy of $x\mapsto x^2\sin(1/x)$ growing from each endpoint, cut each copy at the last point before the midpoint where its derivative vanishes, join the two by a constant, and set the function equal to $0$ on $C$. The resulting function $V$ is differentiable at every point, $V'$ is bounded, and $V'=0$ on $C$; but near each point of $C$ lie endpoints of removed intervals, near which $V'$ oscillates between values close to $-1$ and $1$. So $V'$ is discontinuous at every point of $C$, a set of positive length, and is not Riemann integrable \parencite[for details see][]{Bressoud2008}. Recovering a function from its derivative, the problem that Lebesgue put in the title of his book of 1904, has no general solution in Riemann's theory.

These are three faces of one defect. The class $\lebRiem[a,b]$ is not closed under the limits that analysis takes: pointwise limits (Dirichlet), limits in the mean (the fat Cantor set), and limits of difference quotients (Volterra). Each classical remedy adds a hypothesis that keeps the limit inside the class: uniform convergence, continuity of the limit, integrability of the limit or of the derivative. These hypotheses are not about the problem at hand; they are about the holes in the space.

\section{The move: measure first, then complete}\label{sec:lebesgue:move}

Lebesgue's idea, announced in a note of 1901 and developed in his thesis \parencite{Lebesgue1901,Lebesgue1902},\sidenote{The note was presented to the Paris Academy of Sciences on 29 April 1901; the thesis appeared in the \emph{Annali di Matematica} in 1902.} was to partition the range of a function instead of its domain (\cref{fig:lebesgue:slices}). Let $f$ be a bounded function on $[a,b]$ with values in $[c,d)$, and choose points $c=y_0<y_1<\dots<y_n=d$ whose gaps $y_i-y_{i-1}$ are at most $\delta$. The level sets
\[
  E_i=\{x\in[a,b]:\ y_{i-1}\le f(x)<y_i\}
\]
partition $[a,b]$, and $f$ lies between $\sum_i y_{i-1}\ind_{E_i}$ and $\sum_i y_i\ind_{E_i}$. If each $E_i$ has a length $m(E_i)$, then the integrals of these two functions differ by
\[
  \sum_i(y_i-y_{i-1})\,m(E_i)\le\delta\,(b-a),
\]
however wildly $f$ oscillates. Riemann's brackets differ by the sum of oscillation times length over the intervals of a partition of the domain, which is small only if $f$ oscillates on little of $[a,b]$. Slicing the range removes the oscillation of $f$ from the problem.

\begin{figure}[ht]
\centering
\begin{tikzpicture}[baseline=0pt,x=0.8cm,y=0.8cm,declare function={lebf(\x)=1.5+0.55*sin((1.6*\x-0.3) r)+0.28*sin((4.3*\x+0.9) r);}]
  \foreach \a in {0,0.7,...,4.3}
    {\fill[seamist] (\a,0) rectangle ({\a+0.7},{lebf(\a+0.35)});
     \draw[sea!55,line width=0.3pt] (\a,0) rectangle ({\a+0.7},{lebf(\a+0.35)});}
  \draw[seadeep,line width=0.9pt,domain=0:5.6,samples=160,smooth] plot (\x,{lebf(\x)});
  \draw[muted,->] (0,0) -- (5.9,0);
  \draw[muted,->] (0,0) -- (0,2.55);
  \node[font=\small,text=muted] at (2.8,-0.55) {slice the domain};
\end{tikzpicture}\hspace{0.5cm}%
\begin{tikzpicture}[baseline=0pt,x=0.8cm,y=0.8cm,declare function={lebf(\x)=1.5+0.55*sin((1.6*\x-0.3) r)+0.28*sin((4.3*\x+0.9) r);}]
  \fill[seamist] (0,1.5) rectangle (5.6,1.75);
  \draw[sea!50,line width=0.3pt] (0,1.5) -- (5.6,1.5) (0,1.75) -- (5.6,1.75);
  \foreach \s/\e in {0/0.15,0.59/0.95,1.93/2.05,4.15/4.27,5.27/5.41}
    {\fill[sea!18] (\s,0) rectangle (\e,{min(lebf(\s),lebf(\e))});
     \draw[seadeep,line width=3pt] (\s,0) -- (\e,0);}
  \draw[seadeep,line width=0.9pt,domain=0:5.6,samples=160,smooth] plot (\x,{lebf(\x)});
  \draw[muted,->] (0,0) -- (5.9,0);
  \draw[muted,->] (0,0) -- (0,2.55);
  \node[left,font=\scriptsize] at (0,1.47) {$y_{i-1}$};
  \node[left,font=\scriptsize] at (0,1.80) {$y_{i}$};
  \node[below,font=\scriptsize] at (0.77,-0.02) {$E_i$};
  \node[font=\small,text=muted] at (2.8,-0.55) {slice the range};
\end{tikzpicture}
\caption{Riemann partitions the domain and needs $f$ to oscillate on little of it. Lebesgue partitions the range: the points where $y_{i-1}\le f<y_i$ form the level set $E_i$ (thick segments on the axis), which may be scattered, and only its length is needed.}
\label{fig:lebesgue:slices}
\end{figure}

Lebesgue later explained the difference with money \parencite{Lebesgue1966}. A merchant can count coins and notes in the order in which they come to hand, as Riemann adds values in the order of the points of the domain, or he can sort them by denomination, count each pile, multiply by its value, and add. With finitely many coins the two methods agree; for a sum of infinitely many indivisibles, Lebesgue observed, the difference is decisive. In the integral the piles are the level sets $E_i$, which may be scattered like dust over $[a,b]$. The new object is humble, a set of points on the horizontal axis with one number attached to it, and Grothendieck's advice to stoop low applies almost literally: look down from the graph to the axis. The price is a length for every such set. Measure first, then integrate.

\begin{definition}[Measure spaces]\label[definition]{def:lebesgue:measure}
A \term{$\sigma$-algebra} on a set $X$ is a family $\cM$ of subsets of $X$ that contains $X$ and is closed under complements and countable unions. A \term{measure} on $\cM$ is a function $\mu\colon\cM\to[0,\infty]$ with $\mu(\varnothing)=0$ that is \term{countably additive}: $\mu(\bigcup_kE_k)=\sum_k\mu(E_k)$ for pairwise disjoint $E_k\in\cM$. A set $N\in\cM$ with $\mu(N)=0$ is a \term{null set}, and a property holds \term{almost everywhere} (a.e.) if it fails only on a null set.
\end{definition}

\begin{theorem}[Lebesgue measure]\label{thm:lebesgue:measure}
There are a $\sigma$-algebra $\cL$ of subsets of $\R$ that contains every interval and a measure $m$ on $\cL$ with the following properties.
\begin{enumerate}
\item Length: $m(I)=\abs{I}$ for every interval $I$.
\item Completeness: every subset of a null set belongs to $\cL$.
\item Regularity: for every $E\in\cL$ with $m(E)<\infty$ and every $\eps>0$ there is a finite union $U$ of intervals with $m(E\lebsymdiff U)<\eps$.
\item Invariance: for every $E\in\cL$ and $t\in\R$, the translate $E+t$ belongs to $\cL$ and $m(E+t)=m(E)$.
\end{enumerate}
\end{theorem}
\index{Lebesgue measure}

\begin{proofsketch}
For $A\subseteq\R$ let $m^*(A)$ be the infimum of $\sum_k\abs{I_k}$ over all countable covers of $A$ by open intervals $I_k$. The one inequality with content is $m^*([a,b])\ge b-a$. It holds because every cover of $[a,b]$ by open intervals has a finite subcover, by the Heine--Borel theorem, and for finite covers the inequality is elementary. The criterion of \textcite{Caratheodory1914} then selects the sets $E$ with $m^*(A)=m^*(A\cap E)+m^*(A\setminus E)$ for every $A\subseteq\R$, and a page of set algebra shows that they form a $\sigma$-algebra $\cL$ containing the intervals, on which $m=m^*$ is countably additive. Sets with $m^*(N)=0$ satisfy the criterion, which gives completeness; regularity follows from the definition of $m^*$ as an infimum over covers by intervals, and invariance from the invariance of lengths of intervals. We omit the set algebra; see \textcite[§§1.4--1.5]{Folland1999}.
\end{proofsketch}

\begin{definition}[The integral]\label[definition]{def:lebesgue:integral}
Let $(X,\cM,\mu)$ be a measure space. A function $f\colon X\to[-\infty,\infty]$ is \term{measurable} if $\{f>t\}\in\cM$ for every $t\in\R$. A \term{simple function} is a finite combination $\varphi=\sum_ja_j\ind_{A_j}$ with $A_j\in\cM$. A simple $\varphi\ge0$ can be written with $a_j\ge0$ and pairwise disjoint $A_j$, and then $\int\varphi\dd\mu=\sum_ja_j\mu(A_j)$, with $0\cdot\infty=0$; the value does not depend on the representation. For measurable $f\ge0$,
\[
  \int f\dd\mu=\sup\Bigl\{\int\varphi\dd\mu:\ \varphi\text{ simple},\ 0\le\varphi\le f\Bigr\}.
\]
A measurable $f$ is \term{integrable} if $\int\abs{f}\dd\mu<\infty$, and then, with $f^\pm=\max(\pm f,0)$, its integral is $\int f\dd\mu=\int f^+\dd\mu-\int f^-\dd\mu$; a complex-valued function is measurable, or integrable, when its real and imaginary parts are, and it is integrated through them. For $1\le p<\infty$, the space $L^p(\mu)$ consists of the measurable $f$ with $\norm{f}_p=(\int\abs{f}^p\dd\mu)^{1/p}<\infty$, two functions being identified when they agree almost everywhere; such an $f$ is finite almost everywhere, and we always choose finite-valued representatives. On $[a,b]$ with Lebesgue measure we write $L^p[a,b]$ and $\int_a^bf$.
\end{definition}
\index{measurable function}

We use routine facts freely \parencite[Chapters~2 and~6]{Folland1999}: sums, products, and pointwise limits of measurable functions are measurable, as are suprema and infima of sequences of them; every nonnegative measurable $f$ is the increasing limit of the simple functions $\min\{n,2^{-n}\floor{2^nf}\}$; the integral is monotone, $\abs{\int h\dd\mu}\le\int\abs{h}\dd\mu$, and a nonnegative $f$ has integral $0$ exactly when $f=0$ almost everywhere;\sidenote{Probabilists say \emph{almost surely}; \cref{ch:kolmogorov} builds probability on this chapter's integral.} linearity follows from the simple case and the monotone convergence theorem below, whose proof does not use it; Minkowski's inequality makes $\norm{\cdot}_p$ a norm on $L^p(\mu)$; and for $f,g\in L^2(\mu)$ the product $f\overline{g}$ is integrable, with $\abs{\int f\overline{g}\dd\mu}\le\norm{f}_2\norm{g}_2$.

\subsection{The second face: completion}\label{sec:lebesgue:completionview}

Lebesgue's description is constructive: measure the level sets, then add. A second description says what the construction achieves.

\begin{definition}[Completion]\label[definition]{def:lebesgue:completion}
A normed space $V$ is \term{complete} if every Cauchy sequence in $V$ converges to an element of $V$. A \term{completion} of $V$ is a complete normed space that contains $V$, with its norm, as a dense subspace.
\end{definition}
\index{completion}

Every normed space has a completion, unique up to an isometry that fixes $V$: take the Cauchy sequences in $V$ and identify two of them when their difference tends to zero. The same construction for metric spaces produces $\R$ from $\Q$; Exercise~\ref{ex:lebesgue:reals} shows what goes wrong inside $\Q$. We shall prove in \cref{sec:lebesgue:complete} that $L^1[a,b]$ is complete and that step functions are dense in it. So $L^1[a,b]$ is the completion of the step functions (those that differ at finitely many points being identified), and since $\abs{\int f-\int g}\le\norm{f-g}_1$, the Lebesgue integral is the unique continuous extension of the elementary integral (\cref{cor:lebesgue:completion}). Measuring first and completing the step functions are two descriptions of one object.

The second description makes two points that the first one hides. First, on a complete space where the integral is a continuous linear functional, exchanging a limit with the integral is a matter of continuity once the sequence converges in $\norm{\cdot}_1$; the monotone convergence theorem (when the limit has finite integral) and the dominated convergence theorem below are criteria that turn pointwise convergence into convergence in $L^1$. Second, abstract completion is cheap; what is not cheap is that the completion of the step functions consists of \emph{functions}, every Cauchy sequence of step functions having a subsequence that converges almost everywhere to a representative of its limit. F.~Riesz built the integral on this fact, from two lemmas on monotone sequences of step functions \parencite{RieszSzNagy1952}. The first, in its simplest form, says that step functions decreasing to $0$ at every point have integrals decreasing to $0$. It is the monotone case of Arzelà's theorem for step functions, it needs compactness, as the aside did, and it is countable additivity of length in disguise.

What would one have had to notice to find the move? That the hardest argument of \cref{sec:lebesgue:before}, the aside, manufactures a length for sets that are not intervals, and that the fat Cantor set asks for one, while Riemann's integral knows only the lengths of intervals. A theory that measures more sets integrates more functions, and a theory in which length is countably additive makes the limit theorems automatic.

\section{After: the limit theorems are continuity theorems}\label{sec:lebesgue:limits}

We now repair the first leak. Throughout, $(X,\cM,\mu)$ is a measure space; the reader may keep $[a,b]$ with Lebesgue measure in mind. All the limit theorems in this section draw on one property of $\mu$.

\begin{lemma}[Continuity from below]\label[lemma]{lem:lebesgue:continuity}
If $E_1\subseteq E_2\subseteq\cdots$ are measurable, then $\mu(\bigcup_nE_n)=\lim_n\mu(E_n)$.
\end{lemma}

\begin{proof}
The sets $F_1=E_1$ and $F_n=E_n\setminus E_{n-1}$ for $n\ge2$ are measurable and pairwise disjoint, $E_n=F_1\cup\dots\cup F_n$, and $\bigcup_nE_n=\bigcup_nF_n$. Countable additivity gives $\mu(\bigcup_nE_n)=\sum_k\mu(F_k)=\lim_n\sum_{k\le n}\mu(F_k)=\lim_n\mu(E_n)$.
\end{proof}

\begin{theorem}[Monotone convergence, \textcite{Levi1906}]\label{thm:lebesgue:mct}
Let $0\le f_1\le f_2\le\cdots$ be measurable functions on $X$ with pointwise limit $f$. Then $\int f\dd\mu=\lim_n\int f_n\dd\mu$.
\end{theorem}
\index{monotone convergence theorem}

\begin{proof}
The plan is to test the $f_n$ against a fixed simple function below $f$, on the sets where $f_n$ has nearly caught up with it, and to let continuity from below do the rest. The limit $f$ is measurable, because $\{f>t\}=\bigcup_n\{f_n>t\}$. The integrals $\int f_n\dd\mu$ increase and are at most $\int f\dd\mu$, by monotonicity, so their limit $L$ satisfies $L\le\int f\dd\mu$.

For the reverse inequality, fix a simple function $\varphi=\sum_ja_j\ind_{A_j}$, with $a_j\ge0$ and the $A_j$ pairwise disjoint, such that $0\le\varphi\le f$, and a number $c\in(0,1)$. The sets $E_n=\{f_n\ge c\varphi\}$ are measurable and increase with $n$. Their union is $X$. A point where $\varphi$ vanishes lies in $E_1$, and a point where $\varphi(x)>0$ satisfies $c\varphi(x)<f(x)=\lim_nf_n(x)$, so it lies in $E_n$ for large $n$. The simple function $c\varphi\ind_{E_n}$ lies below $f_n$, so the definition of the integral and \cref{lem:lebesgue:continuity} give
\[
  \int f_n\dd\mu\;\ge\;c\sum_ja_j\,\mu(A_j\cap E_n)\;\longrightarrow\;c\sum_ja_j\,\mu(A_j)=c\int\varphi\dd\mu .
\]
Hence $L\ge c\int\varphi\dd\mu$. Letting $c\to1$ and taking the supremum over $\varphi$ gives $L\ge\int f\dd\mu$.
\end{proof}

The factor $c$ is there for the points where $f_n(x)<f(x)=\varphi(x)$ for every $n$.

\begin{corollary}[Term-by-term integration]\label[corollary]{cor:lebesgue:series}
For every sequence of nonnegative measurable functions $g_k$ on $X$, we have $\int\sum_kg_k\dd\mu=\sum_k\int g_k\dd\mu$.
\end{corollary}

\begin{proof}
The partial sums are measurable, nonnegative, and increase to $\sum_kg_k$. Apply \cref{thm:lebesgue:mct} to them, and additivity of the integral to each partial sum.
\end{proof}

Compare \cref{prop:lebesgue:dini}: no continuity, no compact interval, no Dini argument, and the sum may be infinite. The compactness was paid for once, inside \cref{thm:lebesgue:measure}, and every use of \cref{lem:lebesgue:continuity} draws on that payment.

\begin{theorem}[Fatou's lemma, \textcite{Fatou1906}]\label{thm:lebesgue:fatou}
For every sequence of nonnegative measurable functions $f_n$ on $X$, we have $\int\liminf_nf_n\dd\mu\le\liminf_n\int f_n\dd\mu$.
\end{theorem}
\index{Fatou's lemma}

\begin{proof}
The functions $g_k=\inf_{n\ge k}f_n$ are measurable and increase to $\liminf_nf_n$, and $g_k\le f_n$ for $n\ge k$, so $\int g_k\dd\mu\le\inf_{n\ge k}\int f_n\dd\mu$. Monotone convergence (\cref{thm:lebesgue:mct}) evaluates the left side below, and this bound gives the inequality:
\[
  \int\liminf_nf_n\dd\mu=\lim_k\int g_k\dd\mu\le\lim_k\inf_{n\ge k}\int f_n\dd\mu=\liminf_n\int f_n\dd\mu. \qedhere
\]
\end{proof}

\begin{theorem}[Dominated convergence, \textcite{Lebesgue1908}]\label{thm:lebesgue:dct}
Let $f,f_1,f_2,\dots$ be measurable functions on $X$ with $f_n\to f$ almost everywhere, and suppose that, for every $n$, $\abs{f_n}\le g$ almost everywhere, where $g$ is integrable. Then $f$ is integrable and $\int\abs{f_n-f}\dd\mu\to0$; in particular $\int f_n\dd\mu\to\int f\dd\mu$.
\end{theorem}
\index{dominated convergence theorem}

\begin{proof}
The exceptional sets, together with the null set $\{g=\infty\}$, form a countable family of null sets, and replacing every function by $0$ on their union changes no integral; so we may assume that all the functions are finite and that the convergence and the bounds hold everywhere. Then $\abs{f}\le g$, and $f$ is integrable. The functions $2g-\abs{f_n-f}$ are nonnegative, because $\abs{f_n-f}\le\abs{f_n}+\abs{f}\le2g$, and they converge to $2g$. Fatou's lemma and linearity give
\[
  \int2g\dd\mu\le\liminf_n\int\bigl(2g-\abs{f_n-f}\bigr)\dd\mu=\int2g\dd\mu-\limsup_n\int\abs{f_n-f}\dd\mu .
\]
Since $\int2g\dd\mu$ is finite, $\limsup_n\int\abs{f_n-f}\dd\mu\le0$. The last assertion follows from $\abs{\int f_n\dd\mu-\int f\dd\mu}\le\int\abs{f_n-f}\dd\mu$.
\end{proof}

The theorem says what uniform convergence was standing in for.\sidenote{Lebesgue's thesis has the case of a bounded sequence; the hypothesis $\abs{f_n}\le\abs{F}$ with $F$ summable appears in his paper of 1908.} A convergent sequence can lose mass by concentrating it, as $n\ind_{(0,1/n)}$ does on $[0,1]$, by sending it to infinity, as $\ind_{[n,n+1]}$ does on $\R$, or by spreading it thin, as $\frac1n\ind_{[0,n]}$ does on $\R$; in each example Fatou's inequality is strict (Exercise~\ref{ex:lebesgue:escape}). Domination by an integrable function forbids all three escapes. On $[a,b]$ it is much weaker than uniform convergence, and on $\R$ uniform convergence is not even enough, as the third example shows.

\begin{lemma}[Riemann's integral is a restriction of Lebesgue's]\label[lemma]{lem:lebesgue:riemann}
Every $f\in\lebRiem[a,b]$ is Lebesgue integrable on $[a,b]$, and its Lebesgue integral equals its Riemann integral.
\end{lemma}

\begin{proof}
For step functions this holds because $m(I)=\abs{I}$ for every interval $I$. In general, choose step functions $\varphi_k\le f\le\psi_k$ with $\int(\psi_k-\varphi_k)<1/k$; replacing $\varphi_k$ by $\max(\varphi_1,\dots,\varphi_k)$ and $\psi_k$ by $\min(\psi_1,\dots,\psi_k)$, we may assume that $\varphi_k$ increases and $\psi_k$ decreases. Their limits $\varphi\le f\le\psi$ are measurable, and all these functions are bounded by the constant $\sup\abs{\varphi_1}+\sup\abs{\psi_1}$, which is integrable on $[a,b]$. By \cref{thm:lebesgue:dct}, $\int_a^b(\psi-\varphi)=\lim_k\int(\psi_k-\varphi_k)=0$, so $\psi=\varphi$ almost everywhere, and therefore $f=\varphi$ almost everywhere. Since Lebesgue measure is complete (\cref{thm:lebesgue:measure}), $f$ is measurable, and $\int_a^bf=\int_a^b\varphi=\lim_k\int\varphi_k$, which is the Riemann integral of $f$.
\end{proof}

\begin{corollary}[Arzelà's theorem, in two lines]\label[corollary]{cor:lebesgue:arzela}
Let $f_n\in\lebRiem[a,b]$ satisfy $\abs{f_n}\le M$ for all $n$, and suppose that $f_n\to f$ pointwise. Then $f$ is Lebesgue integrable and $\int_a^bf_n\to\int_a^bf$. When $f\in\lebRiem[a,b]$, this is \cref{thm:lebesgue:arzela}.
\end{corollary}

\begin{proof}
By \cref{lem:lebesgue:riemann}, the $f_n$ are Lebesgue integrable with the same integrals, and their limit $f$ is measurable. The constant $M$ dominates every $\abs{f_n}$ and is integrable on $[a,b]$, so \cref{thm:lebesgue:dct} applies; when $f\in\lebRiem[a,b]$, its Lebesgue integral is its Riemann integral, by \cref{lem:lebesgue:riemann} again.
\end{proof}

Two lines replace the aside, and the hypothesis that the limit is integrable has disappeared. Dirichlet's sequence now behaves as it should: \cref{thm:lebesgue:mct} gives $\int_0^1D=\lim_n\int_0^1f_n=0$.

\section{After: complete spaces, Fourier series, and derivatives}\label{sec:lebesgue:complete}

The second and third leaks close as quickly as the first, beginning with the completeness that Riemann's class lacked. \textcite{Fischer1907} proved it for $L^2$;\sidenote{Fischer had presented the result in a talk at Brünn on 5 March 1907, and in his note he conceded to Riesz the priority of publication.} \textcite{Riesz1907}, independently and a few weeks earlier in the same volume of the \emph{Comptes Rendus}, proved the closely related statement about orthonormal expansions of which \cref{thm:lebesgue:fourier}(i) is the trigonometric case.

\begin{theorem}[Riesz--Fischer]\label{thm:lebesgue:rieszfischer}
For $1\le p<\infty$ the space $L^p(\mu)$ is complete. Moreover, every Cauchy sequence in $L^p(\mu)$ has a subsequence that converges almost everywhere to its limit.
\end{theorem}
\index{Riesz--Fischer theorem}

\begin{proof}
The plan is to pass to a subsequence whose increments are summable, to show with monotone convergence that the telescoping series converges almost everywhere, and to identify its sum as the limit with dominated convergence. Let $(f_n)$ be Cauchy in $L^p(\mu)$, and choose $n_1<n_2<\cdots$ with $\norm{f_{n_{k+1}}-f_{n_k}}_p<2^{-k}$. Put
\[
  G=\abs{f_{n_1}}+\sum_{k\ge1}\abs{f_{n_{k+1}}-f_{n_k}},
\]
a measurable function with values in $[0,\infty]$. By Minkowski's inequality each partial sum $G_K$ of this series satisfies $\norm{G_K}_p\le\norm{f_{n_1}}_p+1$, and the functions $G_K^p$ increase to $G^p$, so \cref{thm:lebesgue:mct} gives $\int G^p\dd\mu\le(\norm{f_{n_1}}_p+1)^p<\infty$. Hence $G$ is finite almost everywhere. Where $G$ is finite, the series $f_{n_1}+\sum_k(f_{n_{k+1}}-f_{n_k})$ converges absolutely; let $f$ be its sum there, which is $\lim_kf_{n_k}$, and let $f=0$ elsewhere. Since $\abs{f_{n_k}}\le G$ and $\abs{f}\le G$, the functions $\abs{f_{n_k}-f}^p$ are dominated by the integrable function $(2G)^p$ and tend to $0$ almost everywhere, so \cref{thm:lebesgue:dct} gives $\norm{f_{n_k}-f}_p\to0$; in particular $f\in L^p(\mu)$. Finally, a Cauchy sequence with a convergent subsequence converges: $\norm{f_n-f}_p\le\norm{f_n-f_{n_k}}_p+\norm{f_{n_k}-f}_p$, and both terms are small when $n$ and $k$ are large.
\end{proof}

\begin{corollary}[$L^1$ is the completion of the step functions]\label[corollary]{cor:lebesgue:completion}
The step functions are dense in $L^1[a,b]$ and in $L^2[a,b]$. Consequently $L^1[a,b]$ is a completion of the space of step functions with the norm $\norm{\cdot}_1$ (step functions that differ at finitely many points being identified), and the Lebesgue integral is the unique continuous linear functional on $L^1[a,b]$ that agrees with the elementary integral on step functions.
\end{corollary}

\begin{proof}
Let $p\in\{1,2\}$ and let $f\in L^p[a,b]$ be real-valued; complex-valued functions are handled through their real and imaginary parts. Applying the construction $\min\{n,2^{-n}\floor{2^ng}\}$ to $g=f^+$ and to $g=f^-$ produces simple functions $\varphi_n\to f$ with $\abs{\varphi_n}\le\abs{f}$, and dominated convergence, with dominating function $(2\abs{f})^p$, gives $\norm{\varphi_n-f}_p\to0$. A simple function is a finite combination of indicators $\ind_E$ with $E\subseteq[a,b]$ measurable, and by the regularity in \cref{thm:lebesgue:measure} there is a finite union $U$ of intervals with $\norm{\ind_E-\ind_U}_p^p=m(E\lebsymdiff U)<\eps$; replacing $U$ by $U\cap[a,b]$ does not increase $m(E\lebsymdiff U)$. So the step functions are dense. They sit in $L^1[a,b]$ isometrically, by \cref{lem:lebesgue:riemann}, and a step function that vanishes almost everywhere vanishes on its open intervals of constancy, so the identification is the one in the statement. Completeness is \cref{thm:lebesgue:rieszfischer}. The integral is linear and satisfies $\abs{\int f-\int g}\le\norm{f-g}_1$, so it is continuous, and two continuous functionals that agree on a dense set agree everywhere.
\end{proof}

\begin{example}[The escaped limit, recovered]\label[example]{ex:lebesgue:cantorlimit}
The step functions $s_n=\ind_{C_n}$ of \cref{prop:lebesgue:incomplete} decrease to $\ind_C$, so dominated convergence gives $m(C)=\lim_n\int s_n=\frac12$ and $\norm{s_n-\ind_C}_1\to0$: the Cauchy sequence that had no limit in $\lebRiem[0,1]$ converges in $L^1$ to the indicator of the fat Cantor set. No Riemann integrable $g$ agrees with $\ind_C$ almost everywhere, for it would satisfy $\norm{s_n-g}_1=\norm{s_n-\ind_C}_1\to0$, by \cref{lem:lebesgue:riemann}, against \cref{prop:lebesgue:incomplete}. The hole in Riemann's class had a definite shape.
\end{example}

\subsection{Fourier series}\label{sec:lebesgue:fourier}

Work in $L^2[0,1]$ with complex-valued functions, the inner product $\ip{f}{g}=\int_0^1f\,\overline{g}$, and the orthonormal functions $e_k(x)=e^{2\pi ikx}$, $k\in\Z$. The Fourier coefficients of $f$ are $\hat f(k)=\ip{f}{e_k}$, and a trigonometric polynomial is a finite linear combination of the $e_k$.

\begin{lemma}[Density of trigonometric polynomials]\label[lemma]{lem:lebesgue:trig}
The trigonometric polynomials are dense in $L^2[0,1]$.
\end{lemma}

\begin{proof}
By \cref{cor:lebesgue:completion} the step functions are dense. For small $\eta>0$, the indicator of $[\alpha,\beta]\subseteq[0,1]$ is within $\sqrt{2\eta}$, in $L^2$, of the continuous function that is $1$ on $[\alpha+\eta,\beta-\eta]$, $0$ outside $(\alpha,\beta)$, and linear in between; so the continuous $g$ with $g(0)=g(1)=0$ are dense. By Fejér's theorem, each continuous $g$ with $g(0)=g(1)$ is a uniform limit of trigonometric polynomials, the averages of the partial sums of its Fourier series \parencite[Chapter~8]{Folland1999}, and uniform convergence on $[0,1]$ implies convergence in $L^2$.
\end{proof}

\begin{theorem}[Riesz--Fischer and Parseval for Fourier series]\label{thm:lebesgue:fourier}
\leavevmode
\begin{enumerate}
\item For every square-summable sequence $(c_k)_{k\in\Z}$, the partial sums $\sum_{\abs{k}\le N}c_ke_k$ converge in $L^2[0,1]$ to a function $f$ with $\hat f(k)=c_k$ for all $k$.
\item For every $f\in L^2[0,1]$, the partial sums $S_Nf=\sum_{\abs{k}\le N}\hat f(k)e_k$ converge to $f$ in $L^2[0,1]$, and $\sum_k\abs{\hat f(k)}^2=\norm{f}_2^2$.
\end{enumerate}
Consequently $f\mapsto(\hat f(k))_{k\in\Z}$ is a linear isometry of $L^2[0,1]$ onto $\ell^2(\Z)$.
\end{theorem}
\index{Parseval's identity}

\begin{proof}
The plan is to let orthonormality make the partial sums Cauchy, completeness supply their limit, and density identify it. For (i), if $M<N$, then orthonormality gives $\norm{\sum_{M<\abs{k}\le N}c_ke_k}_2^2=\sum_{M<\abs{k}\le N}\abs{c_k}^2$, which tends to $0$; so the partial sums are Cauchy and converge to some $f$ by \cref{thm:lebesgue:rieszfischer}. Since the inner product is continuous, $\hat f(j)=\lim_N\ip{\sum_{\abs{k}\le N}c_ke_k}{e_j}=c_j$.

For (ii), the function $f-S_Nf$ is orthogonal to $e_k$ for $\abs{k}\le N$, hence to $S_Nf$, so $\norm{f}_2^2=\norm{S_Nf}_2^2+\norm{f-S_Nf}_2^2$. In particular $\sum_{\abs{k}\le N}\abs{\hat f(k)}^2\le\norm{f}_2^2$ for every $N$, which is Bessel's inequality, and $(\hat f(k))$ is square summable. By (i), $S_Nf$ converges to some $g\in L^2[0,1]$ with $\hat g=\hat f$. The difference $h=f-g$ is orthogonal to every $e_k$, hence to every trigonometric polynomial $P$, and so $\norm{h}_2^2=\ip{h}{h-P}\le\norm{h}_2\norm{h-P}_2$. By \cref{lem:lebesgue:trig} the right side can be made as small as we please, so $h=0$. Thus $S_Nf\to f$, and $\norm{f}_2^2=\lim_N\norm{S_Nf}_2^2=\sum_k\abs{\hat f(k)}^2$. The final assertion combines (i), which gives surjectivity, with (ii), which gives the isometry.
\end{proof}

Part (i) is the theorem that fails in Riemann's class (\cref{sec:lebesgue:incomplete}); its proof uses completeness, orthogonality, and nothing else. Fourier analysis on $[0,1]$ has become coordinate geometry in a Hilbert space.

\subsection{The fundamental theorem for bounded derivatives}\label{sec:lebesgue:ftc}

\begin{theorem}[\textcite{Lebesgue1902}]\label{thm:lebesgue:ftc}
Let $F\colon[a,b]\to\R$ be differentiable at every point of $[a,b]$, one-sidedly at the endpoints, and suppose that $\abs{F'}\le M$ on $[a,b]$. Then $F'$ is Lebesgue integrable and $\int_a^bF'=F(b)-F(a)$.
\end{theorem}
\index{fundamental theorem of calculus}

\begin{proof}
The plan is to write $F'$ as a pointwise limit of difference quotients, which are continuous, bounded by the mean value theorem, and integrable by hand, and then to apply dominated convergence. Extend $F$ to $[a,b+1]$ by $F(x)=F(b)+F'(b)(x-b)$ for $x>b$. The extension is differentiable on $[a,b+1)$, since at $b$ both one-sided derivatives equal $F'(b)$, and its derivative is bounded by $M$. For $n\ge1$ and $x\in[a,b]$ put $g_n(x)=n\bigl(F(x+\frac1n)-F(x)\bigr)$. Each $g_n$ is continuous, and $g_n(x)\to F'(x)$ for every $x\in[a,b]$ by the definition of the derivative. By the mean value theorem, $g_n(x)=F'(\xi)$ for some $\xi\in(x,x+\frac1n)$, so $\abs{g_n}\le M$. Hence $F'$ is measurable, as a pointwise limit of continuous functions, and \cref{thm:lebesgue:dct}, with the constant $M$ as dominating function, gives $\int_a^bF'=\lim_n\int_a^bg_n$. To compute $\int_a^bg_n$ for $n>1/(b-a)$, shift the variable in the first term; the integrals over $[a+\frac1n,b]$ cancel, and
\[
  \int_a^bg_n=n\int_b^{b+1/n}F-n\int_a^{a+1/n}F\;\longrightarrow\;F(b)-F(a),
\]
because $F$ is continuous.
\end{proof}

Volterra's function is now harmless:\sidenote{Lebesgue's thesis reproduces Volterra's example and credits him with it.} its derivative is bounded, so $\int_0^1V'=V(1)-V(0)$. Lebesgue announced \cref{thm:lebesgue:ftc} in his note of 1901 and proved it in his thesis, by the argument above with his bounded convergence theorem in place of dominated convergence. Boundedness cannot simply be dropped (Exercise~\ref{ex:lebesgue:unbounded}), nor can differentiability at every point be weakened to differentiability almost everywhere, as the Cantor function of \cref{sec:lebesgue:costs} shows.

\section{Costs and limits}\label{sec:lebesgue:costs}

\paragraph{Where did the difficulty go?} Into one inequality and one page of set algebra. The inequality is $m^*([a,b])\ge b-a$: a countable cover of an interval by open intervals has total length at least the length of the interval. It fails for $\Q\cap[a,b]$, which intervals of arbitrarily small total length can cover, and it holds for $[a,b]$ because $[a,b]$ is compact. Carathéodory's set algebra turns it into a countably additive measure, and every limit theorem of \cref{sec:lebesgue:limits} draws on this single investment through \cref{lem:lebesgue:continuity}; Riemann's world paid for the same compactness again and again, in \cref{prop:lebesgue:dini} and inside an induction in the aside. Paying once is the bargain. A quieter cost is measurability: functions must be shown measurable before they are integrated, and an element of $L^1[a,b]$ is an equivalence class, whose value at a single point means nothing. The completion bought its limits by giving up points.

\paragraph{What did not dissolve.} Four things. First, not every set has a length.\index{non-measurable set} Using the axiom of choice, \textcite{Vitali1905} picked one point in $[0,1]$ from each coset of $\Q$ in $\R$. The translates of the resulting set $W$ by the rationals in $[-1,1]$ are pairwise disjoint, cover $[0,1]$, and lie in $[-1,2]$, so countable additivity and invariance leave no possible value for $m(W)$: zero would give $m([0,1])=0$, and a positive value would give $m([-1,2])=\infty$. Some choice is unavoidable here: \textcite{Solovay1970} showed that if set theory with an inaccessible cardinal is consistent, then so is set theory with only the axiom of dependent choice in which every set of reals is Lebesgue measurable, and Shelah later showed that the inaccessible cannot be dropped.\sidenote{S.~Shelah, ``Can you take Solovay's inaccessible away?'', \emph{Israel Journal of Mathematics} 48 (1984).} Second, the integral is absolute: $f$ is integrable only if $\abs{f}$ is. The improper Riemann integral of $\frac{\sin x}{x}$ over $(0,\infty)$ is $\frac\pi2$, but $\int_0^\infty\frac{\abs{\sin x}}{x}\dd x=\infty$, so $\frac{\sin x}{x}$ is not Lebesgue integrable there. Third, not every derivative is integrable: $F(x)=x^2\sin(x^{-2})$ is differentiable at every point of $[0,1]$, but $F'$ is not Lebesgue integrable (Exercise~\ref{ex:lebesgue:unbounded}). The integral of Denjoy and Perron, which Kurzweil and Henstock later recast as a limit of Riemann sums,\sidenote{Kurzweil (1957) arrived at the definition while studying differential equations; Henstock (1961) found it independently.} integrates every derivative \parencite{Bartle2001},\index{Henstock--Kurzweil integral} but its natural normed space, under the Alexiewicz norm $\sup_x\abs{\int_a^xf}$, is not complete. So ``complete in which norm?'' is a real choice: Lebesgue chose completeness, and the gauge integral chose every derivative. Fourth, the general fundamental theorem, for absolutely continuous $F$, needs a covering lemma that completion does not supply; the Cantor function, which increases from $0$ to $1$ with derivative $0$ almost everywhere, marks its boundary.

\paragraph{Other seas.} \textcite{Daniell1918} skipped measure: extend a positive linear functional $I$ from a vector lattice of elementary functions to monotone limits, and recover measures afterwards as integrals of indicators. The extension needs one more hypothesis, $I(h_n)\to0$ whenever $h_n$ decreases to $0$, which is Riesz's first lemma again, the place where compactness is paid. This is the completion picture taken as a definition. Bochner's integral runs the same construction for functions with values in a Banach space, and Schwartz's distributions complete in another direction, until differentiation rather than integration is continuous. The operation that one decides to make continuous decides the completion.

\section{The move, distilled}\label{sec:lebesgue:distilled}

\begin{move}{Completion}
\movepart{In one sentence} When limits keep escaping your space, put the limits into the space: complete it in a norm in which the operations you care about are continuous.
\movepart{Triggers} Theorems that need uniform convergence, or a hypothesis such as ``assume the limit is integrable'', to interchange two limits; natural Cauchy sequences with no limit; existence proofs that need a minimizing sequence to converge; hypotheses that only someone who knows the answer can check.
\movepart{Recipe} Find the norm in which the operations you care about are continuous (here, $\norm{\cdot}_1$). Complete. Identify the new points concretely (here, as almost-everywhere limits, that is, measurable functions) and pay for that identification once. Check that the old objects are dense. Prove limit theorems as continuity statements, and extend the rest by density.
\movepart{Anti-pattern} Completing in a norm in which the operations you care about are not continuous, or the sequences you care about do not converge (Exercise~\ref{ex:lebesgue:graphs}); or completing abstractly, with no concrete description of the new points, so that every application must reconstruct them.
\end{move}

\dissolutiontest{Arzelà's theorem: from an inner-content lemma with a compactness argument to two lines. Term-by-term integration: from Dini's theorem to one line.}{Domination replaces uniform convergence; the limit need not be continuous or Riemann integrable; every bounded derivative integrates to $F(b)-F(a)$.}{Every interchange of a limit and an integral is continuity of one linear functional on one complete space; Fourier series are coordinates in a Hilbert space.}{The Riesz--Fischer theorem, $L^2\cong\ell^2$; the fundamental theorem for bounded derivatives; probability as integration (\cref{ch:kolmogorov}).}

Completion recurs in this book. \Cref{ch:graphon} completes the space of graphs in the cut norm, in which subgraph densities are continuous, and so avoids the anti-pattern of Exercise~\ref{ex:lebesgue:graphs}; \cref{ch:transport} works with probability measures in the weak topology, where optimal couplings exist by compactness (Exercise~\ref{ex:lebesgue:coupling}); the pseudo-distributions of \cref{ch:sos} add ideal points to a space of distributions. \Cref{ch:kolmogorov} takes this chapter's integral as given and calls it expectation; its move is axiomatization rather than completion, but its strong law speaks of almost-everywhere limits, which only the completed theory can express.

\begin{exercises}
\item\label{ex:lebesgue:escape} \consolidation (Escaping mass.) (a) On $[0,1]$ let $f_n=n\ind_{(0,1/n)}$. Compute $\int\lim_nf_n$ and $\lim_n\int f_n$, and show that $\sup_nf_n$ is not integrable. (b) Do the same on $\R$ for $f_n=\ind_{[n,n+1]}$ and for $f_n=\frac1n\ind_{[0,n]}$, which converges uniformly. (c) For $k\ge0$ and $0\le j<2^k$ let $h_{k,j}$ be the indicator of $[j2^{-k},(j+1)2^{-k}]$, and list these functions, ordered by $k$ and then by $j$, as $g_1,g_2,\dots$. Show that $\norm{g_n}_1\to0$ but that $(g_n(x))$ converges for no $x\in[0,1]$, and find a subsequence that converges to $0$ almost everywhere, as \cref{thm:lebesgue:rieszfischer} promises.

\item\label{ex:lebesgue:fourier} \consolidation (Coefficients without a Riemann function.) Let $C$ be the fat Cantor set of \cref{sec:lebesgue:incomplete} and $c_k=\int_Ce^{-2\pi ikx}\dd x$. Show that $(c_k)$ is square summable, that $c_k=\lim_n\hat s_n(k)$, and that $(c_k)$ is the sequence of Fourier coefficients of no Riemann integrable function on $[0,1]$. \hint{if $\hat g=\hat\ind_C$, then $g=\ind_C$ almost everywhere by \cref{thm:lebesgue:fourier}; now use \cref{ex:lebesgue:cantorlimit}.}

\item\label{ex:lebesgue:lp} \consolidation (Sequence spaces.) For counting measure $\mu$ on $\N$, show that $L^p(\mu)=\ell^p$ and that ``almost everywhere'' means ``everywhere''; deduce the completeness of $\ell^p$ from \cref{thm:lebesgue:rieszfischer}. Then prove it directly, and say which step of the general proof becomes invisible, and why.

\item\label{ex:lebesgue:unbounded} \consolidation (An unbounded derivative.) Let $F(x)=x^2\sin(x^{-2})$ for $0<x\le1$ and $F(0)=0$. (a) Show that $F$ is differentiable at every point of $[0,1]$. (b) Show that $F'$ is not Lebesgue integrable on $[0,1]$. \hint{substitute $u=x^{-2}$ in $\int_0^1x^{-1}\abs{\cos(x^{-2})}\dd x$.} (c) Show that $\lim_{\eps\to0}\int_\eps^1F'=F(1)-F(0)$, so that the fundamental theorem survives as an improper integral.

\item\label{ex:lebesgue:reals} \transfer (The real numbers as a completion.) (a) Show that the rationals $x_0=1$, $x_{n+1}=\frac12(x_n+2/x_n)$ form a Cauchy sequence with no rational limit. (b) The polynomial $x^2-2$ changes sign on $[1,2]$ but has no rational zero: the intermediate value theorem leaks from $\Q$. Prove it on $\R$ by bisection, and point to the step that uses completeness. (c) Name two more theorems of calculus that fail over $\Q$. Where did the difficulty go?

\item\label{ex:lebesgue:dirichlet} \transfer\hard (The Dirichlet principle.) Let $V$ be the space of continuously differentiable $u$ on $[0,1]$ with $u(0)=u(1)=0$, with the inner product $\ip{u}{v}=\int_0^1u'v'$, and let $f$ be continuous on $[0,1]$. (a) Prove that $\abs{u(x)}\le\norm{u}$ for $u\in V$ and $x\in[0,1]$, and show that $V$ is not complete. (b) Show that every element of the completion $H$ of $V$ is realized by a continuous function, the uniform limit of a Cauchy sequence in $V$, and that distinct elements give distinct functions. (c) Using that $H$ is a Hilbert space and that $u\mapsto\int_0^1fu$ is continuous on $H$, by (a), show that $J(u)=\frac12\norm{u}^2-\int_0^1fu$ has exactly one minimizer in $H$. (d) Check that the minimizer is $u(x)=\int_0^1\min(x,y)\bigl(1-\max(x,y)\bigr)f(y)\dd y$, which solves $-u''=f$. In higher dimensions (c) is just as easy and (d) becomes a regularity theorem: say where the difficulty went. \hint{for (b), use the completeness of $L^2[0,1]$ and the density of step functions.}

\item\label{ex:lebesgue:coupling} \transfer (Existence of optimal couplings.) For Borel probability measures on a compact metric space $K$, say that $\pi_n\to\pi$ weakly if $\int\phi\dd\pi_n\to\int\phi\dd\pi$ for every continuous $\phi$. (a) Show that $\delta_{1/n}\to\delta_0$ weakly on $[0,1]$, although $\abs{\delta_{1/n}(A)-\delta_0(A)}=1$ for $A=\{0\}$: in total variation, Dirac masses never approach each other. (b) Prokhorov's theorem says that every sequence of Borel probability measures on $K$ has a weakly convergent subsequence. Use it to show that for Borel probability measures $\mu,\nu$ on $[0,1]$ and continuous $c$ on $[0,1]^2$, some coupling of $\mu$ and $\nu$, that is, some probability measure $\pi$ on $[0,1]^2$ with marginals $\mu$ and $\nu$, minimizes $\int c\dd\pi$. You may use that the product measure $\mu\otimes\nu$ is a coupling, and that a Borel probability measure on $K$ is determined by the integrals of continuous functions. Optimal couplings are the subject of \cref{ch:transport}.

\item\label{ex:lebesgue:graphs} \transfer (Graphs in the wrong norm.) Divide $[0,1]$ into intervals $I_1,\dots,I_n$ of length $\frac1n$, and for a graph $G$ on $\{1,\dots,n\}$ let $W_G$ be $1$ on $I_i\times I_j$ when $ij$ is an edge and $0$ elsewhere. (a) Show that $\norm{W_G-\frac12}_{L^1}=\frac12$ for every graph $G$. (b) For measurable $U,W\colon[0,1]^2\to[0,1]$ let $t(W)=\int W(x,y)W(y,z)W(z,x)\dd x\dd y\dd z$; show, with Fubini's theorem, that $\abs{t(U)-t(W)}\le3\norm{U-W}_{L^1}$. (c) For a uniformly random graph $G_n$ on $\{1,\dots,n\}$, show that $\E\,t(W_{G_n})\to\frac18=t(\frac12)$. Why is $L^1$ the wrong norm in which to complete the space of graphs? Compare the cut norm of \cref{ch:graphon}.

\item\label{ex:lebesgue:meanfield} \research\harder (Completing neural networks.) A two-layer network $f(x)=\frac1N\sum_{i\le N}a_i\sigma(\ip{w_i}{x})$ integrates $a\,\sigma(\ip{w}{x})$ against the empirical measure of its parameters $(a_i,w_i)$. Letting this measure range over all probability measures completes the finite networks, and as $N\to\infty$ suitably scaled gradient descent becomes a Wasserstein gradient flow of the risk \parencite{MeiMontanariNguyen2018,ChizatBach2018}. (a) For bounded continuous $\sigma$, data with $\E y^2<\infty$, and measures on a fixed compact set, show that $R(\mu)=\E\bigl(y-\int a\,\sigma(\ip{w}{x})\dd\mu(a,w)\bigr)^2$ is continuous under weak convergence. (b) Open: for deep networks, which completion makes both the network function and the training dynamics continuous, and which global-convergence guarantees known for two layers survive in it?
\end{exercises}

\notesandrefs

\textcite{Hawkins1970} is the standard history of how Lebesgue's theory grew out of the problems of \cref{sec:lebesgue:before}, and \textcite{Bressoud2008} builds a course around them, with Volterra's construction in detail. Lebesgue's note of 1901 and thesis of 1902 contain the measure, the integral, the bounded convergence theorem, and the integration of bounded derivatives; his \emph{Leçons} \parencite{Lebesgue1904} gave a book-length account.

The monotone convergence theorem for Lebesgue's integral is due to \textcite{Levi1906}, and Fatou's lemma comes from Fatou's thesis \parencite{Fatou1906}; \textcite{Folland1999} covers \cref{sec:lebesgue:limits,sec:lebesgue:complete} in modern form. For Arzelà's theorem without measure theory see \textcite{Luxemburg1971}; the aside follows \textcite{Lewin1986}. For the gauge integral see \textcite{Bartle2001}.

\printchapterbib
\end{refsection}
